% paper/sections/12e_interface_crossing.tex
%
% §11.5  高密度比対応：一括解法の限界
% ═══════════════════════════════════════════════════════════════════════════════
\subsection{高密度比対応：変密度一括解法の適用限界}
\label{sec:interface_crossing}
% ═══════════════════════════════════════════════════════════════════════════════

§~\ref{sec:split_ppe_framework} で分相 PPE 解法を主要 PPE 戦略として定式化し，
§~\ref{sec:verify_split_ppe_dc} でその密度比非依存な精度（$\Ord{h^{7.0}}$）を実証した．
本節では，変密度 smoothed Heaviside 一括解法の\textbf{破綻閾値}を定量的に記録し，
分相 PPE 解法の必要性を裏付ける補足データを示す．

% ─────────────────────────────────────────────────────────────────────────────
\subsubsection{界面型密度場での DC 収束特性}
\label{sec:varrho_dc_convergence}
% ─────────────────────────────────────────────────────────────────────────────

\noindent\textbf{目的：}
変密度 smoothed Heaviside 一括解法の欠陥補正（DC）が
界面型密度場で収束可能な密度比の閾値を定量化する．

\noindent\textbf{評価対象：}
変密度 PPE 欠陥補正（DC $k = 3$）の密度比依存性．

\noindent\textbf{手順とパラメタ：}
\begin{itemize}[noitemsep]
  \item 製造解 $p^* = \sin\pi x\sin\pi y$，ディリクレ BC
  \item 密度場 $\rho = \rho_g + (\rho_l - \rho_g)\,H_\varepsilon(\phi)$（$\varepsilon = 1.5h$，円形界面 $R = 0.25$）
  \item $N = 64$，DC $k = 3$
  \item 密度比を $\rho_l/\rho_g = 2, 5, 10, 20, 50, 100$ で変化
\end{itemize}

\noindent\textbf{結果：}

\begin{table}[htbp]
\centering
\caption{変密度 PPE 欠陥補正の密度比依存性（$N = 64$，$k = 3$，界面型密度場）}
\label{tab:varrho_breakdown}
\renewcommand{\arraystretch}{1.3}
\begin{tabular}{rcccl}
\toprule
$\rho_l/\rho_g$ & $L^\infty$ 誤差 & $\|r_0\|_\infty$ & $\|r_3\|_\infty$ & 判定 \\
\midrule
   2 & $1.1 \times 10^{-3}$ & $3.6 \times 10^{1}$ & $1.3 \times 10^{1}$ & 収束（遅化） \\
   5 & $3.8 \times 10^{-2}$ & $8.7 \times 10^{1}$ & $3.6 \times 10^{1}$ & 収束（遅化） \\
  10 & $1.3 \times 10^{0}$  & $1.3 \times 10^{2}$ & $3.4 \times 10^{2}$ & \textbf{発散} \\
  20 & $9.7 \times 10^{1}$  & $1.7 \times 10^{2}$ & $1.1 \times 10^{4}$ & \textbf{発散} \\
  50 & $2.9 \times 10^{2}$  & $2.3 \times 10^{2}$ & $7.3 \times 10^{4}$ & \textbf{発散} \\
 100 & $1.1 \times 10^{5}$  & $2.9 \times 10^{2}$ & $3.6 \times 10^{7}$ & \textbf{発散} \\
\bottomrule
\end{tabular}
\end{table}

\noindent\textbf{結論：}
$\rho_l/\rho_g \geq 10$ では DC が発散し（$\|r_3\| > \|r_0\|$），
反復の増大や緩和係数では回復不能である．
$\rho_l/\rho_g = 2, 5$ でも残差比は $\sim 0.35$--$0.41$ と，
滑らかな密度場（§~\ref{sec:verify_ppe_varrho}，残差比 $\sim 10^{-4}$）比で 3 桁以上劣化している．

% ─────────────────────────────────────────────────────────────────────────────
\subsubsection{格子収束性の密度比依存}
\label{sec:varrho_grid_convergence}
% ─────────────────────────────────────────────────────────────────────────────

\noindent\textbf{目的：}
DC が収束する低密度比（$\rho_l/\rho_g = 2, 5$）について格子収束次数を計測し，
CCD の設計精度 $\Ord{h^6}$ が界面型密度場で回復可能かを検証する．

\noindent\textbf{評価対象：}
格子収束次数の密度比依存性（変密度一括 PPE + CCD 離散化）．

\noindent\textbf{手順とパラメタ：}
\begin{itemize}[noitemsep]
  \item 前節と同一の製造解（$p^* = \sin\pi x\sin\pi y$，界面型密度場）
  \item $\rho_l/\rho_g = 2, 5$（DC 収束領域），格子精緻化
  \item 静止液滴密度比スイープ：$\rho_l/\rho_g = 2, 3, 5, 10$，$N = 64$，$\mathit{We} = 10$，200 ステップ
\end{itemize}

\noindent\textbf{結果：}

$\rho_l/\rho_g = 2$ で $\Ord{h^{1}}$--$\Ord{h^{2}}$，
$\rho_l/\rho_g = 5$ で $\Ord{h^{0.1}}$--$\Ord{h^{0.3}}$ に劣化した．
CCD の設計精度 $\Ord{h^6}$ は界面近傍で回復不能である．

図~\ref{fig:ch12_density_fields} に静止液滴の圧力場と寄生電流を
$\rho_l/\rho_g = 2, 3, 5, 10$ で比較する．
$\rho_l/\rho_g \leq 10$ では Laplace 圧が正しく維持され（$\Delta p_{\mathrm{meas}} = 0.395$，
厳密値 $0.400$），寄生流れは $\Ord{10^{-4}}$ で有界である．
変密度一括解法の適用範囲内であれば，
密度比の増大は安定性に大きな影響を与えない．

\begin{figure}[htbp]
\centering
\includegraphics[width=\linewidth]{figures/ch12_density_fields}
\caption{密度比スイープ：静止液滴の圧力場（上段）と寄生電流（下段）．
$N = 64$，$\mathit{We} = 10$，200 ステップ．
$\rho_l/\rho_g = 2, 3, 5, 10$ の全ケースで安定．}
\label{fig:ch12_density_fields}
\end{figure}

\begin{tcolorbox}[resultbox, title={Smoothed Heaviside 一括解法の適用限界}]
\begin{enumerate}[noitemsep]
  \item $\rho_l/\rho_g \leq 5$：DC 収束するが，
        格子収束次数は $\Ord{h^{0.1}}$--$\Ord{h^{2}}$ に劣化
  \item $\rho_l/\rho_g \geq 10$：DC 発散（反復行列のスペクトル半径 $> 1$）
  \item 水--空気系（$\rho_l/\rho_g \approx 830$）：使用不可
\end{enumerate}
分相 PPE + DC $k{=}3$（§~\ref{sec:split_ppe_framework}）は
全密度比で $\Ord{h^{7.0}}$ を達成し，
この制約を完全に回避する（§~\ref{sec:verify_split_ppe_dc}）．
ただし分相 PPE は\textbf{相ごとに定密度 Poisson 方程式を個別に解く}手法であり，
smoothed Heaviside の変密度演算子を定密度に置換するだけでは
界面での $(\bnabla\rho \cdot \bnabla p)/\rho^2$ 項の欠落により
速度発散誤差が蓄積する．
実装には GFM 等の界面ジャンプ処理が不可欠である．
\end{tcolorbox}

\noindent\textbf{結論：}
格子収束次数は $\rho_l/\rho_g = 2$ で $\Ord{h^{1}}$--$\Ord{h^{2}}$，
$\rho_l/\rho_g = 5$ で $\Ord{h^{0.1}}$--$\Ord{h^{0.3}}$ であり，
CCD の設計精度 $\Ord{h^6}$ は界面近傍で回復不能である．
DC にはもう一つの本質的制約がある：
CCD 高次演算子 $L_H$ と低次前処理 $L_L$ の
スペクトル不整合により定常反復が発散する（密度比非依存；
付録~\ref{sec:ppe_discretization_choice} にスペクトル解析と定量比較を収録）．
FD 直接法 PPE + CCD 勾配の併用が最適である．
